\section{La unidad de reporte}
\label{sec:unidad-reporte}

\subsection{Dos niveles, no uno}

El reporte tiene dos niveles y conviene no mezclarlos, porque tienen funciones
distintas.

El \textbf{panel} es la partición obligatoria: category por aspect, nueve celdas,
siempre las nueve, nunca agregadas. Su función es impedir que un promedio esconda
que el método funciona en un problema y no en otro. Las nueve celdas son nueve
problemas distintos, no nueve muestras del mismo.

El \textbf{estrato} es el refinamiento dentro del panel: los siete ejes de
D\ref{def:estrato}. Su función es distinta, y es la que interesa al proyecto:
describir a cada proteína \emph{por su posición en el vecindario}, no solo por su
conocimiento previo. Una proteína corta con un donante casi idéntico y una
proteína larga sin ningún donante anotado son dos problemas de recuperación
completamente distintos que el panel por sí solo no separa.

El mecanismo que lo hace posible sin violar la regla de no agregación es que el
estrato \textbf{viaja dentro de la clave del panel}. Una celda se identifica como
\cod{NK:MFO@length=512-1024}, de modo que una restricción por estrato produce
nueve celdas restringidas y nunca un número.

\begin{aviso}{La pertenencia al estrato se toma del brazo base, no del comparado}
Si cada arm decidiera la banda de sus propias queries, dos arms estarían
comparándose sobre poblaciones distintas y la diferencia medida incluiría el
cambio de población. La pertenencia se resuelve una vez, sobre el brazo base, y se
aplica a los dos lados.
\end{aviso}

\subsection{El límite real: granularidad contra población}

Aquí está la tensión que el diseño tiene que resolver, y es aritmética.

El cruce completo de los siete ejes tiene
$3 \times 3 \times 4 \times 5 \times 3 \times 6 \times 5 = 16{,}200$ celdas. La
campaña dispone de unas 30{,}427 unidades emparejadas en total. Es decir, el cruce
completo promedia menos de dos proteínas por celda, y el suelo de población para
que una celda tenga voto es 332 con el $\sigma$ global, y entre 120 y 1{,}376 con el
de cada panel (D\ref{def:floorpop}). \textbf{El cruce de
siete ejes está vacío en casi todas partes}, y no por falta de cómputo: por falta
de proteínas.

Medido sobre la campaña actual, cruzando el panel con una sola banda de longitud,
sobreviven \textbf{17 celdas} bajo el suelo global de 332, y solo \textbf{11} bajo
el suelo por panel. Y las once viven en cinco paneles: cuatro en \cod{PK:BPO},
tres en \cod{PK:MFO}, dos en \cod{PK:CCO}, una en \cod{LK:BPO} y una en
\cod{NK:BPO}. \textbf{Cuatro paneles no aportan ninguna región}: \cod{NK:MFO},
\cod{NK:CCO}, \cod{LK:MFO} y \cod{LK:CCO}. Dicho de otro modo, encaminar por
longitud sobre un suelo defendible solo puede encaminar dentro de \cod{PK} y de
\cod{BPO}: no puede encaminar \cod{MFO} ni \cod{CCO} fuera de \cod{PK}.

De esto se siguen tres reglas de diseño:

\begin{enumerate}[leftmargin=1.5em,itemsep=0.35em]
  \item \textbf{El panel es obligatorio; el refinamiento es de un eje a la vez.}
        Panel por longitud, o panel por homología, pero no panel por longitud por
        homología, que no tiene población.
  \item \textbf{El eje del refinamiento se declara antes de la corrida.} Si se
        prueban varios y se reporta el que separó, eso es seleccionar sobre ruido
        con otro nombre. La sección~\ref{sec:estadistico} cuantifica cuánto ruido.
  \item \textbf{El cruce completo es una herramienta de descripción del campeón,
        no de selección.} Una vez elegida una configuración, describirla sobre los
        siete ejes es informativo y no decide nada. Usarla para elegir es lo que no
        se puede hacer.
\end{enumerate}

\subsection{Los dos ejes que el proyecto quiere, y su estado}

\begin{center}
\begin{tabularx}{\linewidth}{@{}L{2.3cm} L{3.4cm} X@{}}
\toprule
eje & precondición & estado\\
\midrule
\cod{length}    & ninguna. Es una propiedad de la query & disponible sobre todas las queries. Cuatro bandas por los límites de contexto\\
\cod{homology}  & \cod{compute\_alignments} activo y el backfill de identidad completo & disponible. Medido tres veces por panel y solo bajo dos de los sellos, lo que es poco\\
\bottomrule
\end{tabularx}
\end{center}

La banda de homología es la que más información aporta sobre el método, porque es
la que separa el régimen donde la transferencia por similitud es casi una copia del
régimen donde deja de implicar función compartida. Un resultado ya medido en la
campaña lo ilustra: la política sin filtro, que es la más permisiva de las dos,
gana en \textbf{196 de las 197 lecturas que resuelven} (la excepción es
\cod{LK:MFO} en la banda \cod{>90}), y gana con un efecto entre un 69 y un 78\,\%
mayor por debajo del 30\,\% de identidad que por encima del 90\,\%, según cómo se
agregue. La banda \cod{none} \textbf{no está medida}: sus seis comparaciones
fallaron, de modo que la homología está medida sobre cuatro bandas y no cinco. El mismo veredicto, con una lectura mucho más informativa.

\subsection{Población y poder por panel}

Toda publicación de un veredicto lleva al lado esta tabla, porque el poder de
resolución de los nueve paneles abarca casi un orden de magnitud y una celda que
no puede resolver nada no debe leerse como una celda que no encontró nada.

\begin{center}
\small
\begin{tabularx}{\linewidth}{@{}l R{1.7cm} R{1.9cm} R{2.5cm} R{1.5cm} R{1.5cm} X@{}}
\toprule
panel & $n$ & semiamp. IC & MDE, rango & $\sigma$ impl. & suelo & vota\\
\midrule
\cod{PK:BPO} & 13{,}876 & 0.00130 & 0.0001 - 0.0041 & 0.0782 &   120 & sí\\
\cod{PK:MFO} &  4{,}947 & 0.00421 & 0.0005 - 0.0121 & 0.1511 &   448 & sí\\
\cod{PK:CCO} &  4{,}872 & 0.00466 & 0.0004 - 0.0127 & 0.1659 &   541 & sí\\
\cod{NK:BPO} &  1{,}509 & 0.00951 & 0.0010 - 0.0222 & 0.1884 &   697 & sí\\
\cod{LK:BPO} &  1{,}214 & 0.01102 & 0.0011 - 0.0292 & 0.1958 &   753 & sí\\
\cod{NK:MFO} &  1{,}129 & 0.01544 & 0.0018 - 0.0339 & 0.2647 & 1{,}376 & \textbf{no}\\
\cod{NK:CCO} &  1{,}116 & 0.01520 & 0.0014 - 0.0413 & 0.2591 & 1{,}317 & \textbf{no}\\
\cod{LK:MFO} &     943  & 0.01375 & 0.0016 - 0.0348 & 0.2154 &   911 & sí\\
\cod{LK:CCO} &     821  & 0.01735 & 0.0016 - 0.0431 & 0.2536 & 1{,}263 & \textbf{no}\\
\bottomrule
\end{tabularx}
\captionof{table}{Poder de resolución por panel, sobre la variante
\cod{b7cfed9a}. Todas las columnas salen del registro: $n$ y la semiamplitud del
intervalo bootstrap de los eventos \cod{pairing} y \cod{panel}, el rango del MDE
del propio campo \cod{minimum\_detectable\_effect} sobre las lecturas que
resuelven, y $\sigma$ implícita despejada de la semiamplitud. El MDE se da como
rango porque no es una constante del panel: hay uno por comparación.}
\end{center}

\begin{aviso}{Tres de los nueve paneles no alcanzan su propio suelo}
La columna de suelo aplica la misma regla que el suelo global,
$\lceil (2.8016\,\sigma/0.02)^2 \rceil$, pero con el $\sigma$ implícito de cada
panel en vez de un 0.13 único. Bajo ese criterio \cod{NK:MFO}, \cod{NK:CCO} y
\cod{LK:CCO} \textbf{no pueden votar}: su población está por debajo del mínimo que
su propia dispersión exige. Dos de los tres son \cod{NK}, que es la categoría que
el proyecto quiere titular.

No es un defecto de la tabla, es lo que la tabla mide. Un suelo global de 332
dejaba votar a las nueve porque suponía una dispersión que ocho de los nueve
paneles no tienen.
\end{aviso}

La consecuencia que gobierna la regla de cierre está en esta tabla. Las cinco
celdas más pequeñas son \emph{mayoría de las nueve} y sostienen solo el 17.2\,\%
de la evidencia. La razón entre el error típico mayor y el menor es 13.4. Una regla
que sellase por mayoría simple de los nueve paneles podría cerrar un eje, convertirse
en el umbral del siguiente y propagarse aguas abajo con el voto de la sexta parte
peor resuelta del registro.
